\chapter{API REST}

A API FastAPI expõe o catálogo em JSON e permite filtros combináveis. A documentação interativa normalmente fica disponível em \texttt{/docs}.

\section{Endpoints principais}

\begin{center}
\small
\renewcommand{\arraystretch}{1.15}
\begin{longtable}{>{\RaggedRight\arraybackslash}p{0.29\textwidth}>{\RaggedRight\arraybackslash}p{0.64\textwidth}}
\toprule
\textbf{\textbf{Método e rota}} & \textbf{\textbf{Finalidade}} \\
\midrule
\endfirsthead
\toprule
\textbf{\textbf{Método e rota}} & \textbf{\textbf{Finalidade}} \\
\midrule
\endhead
\midrule
\multicolumn{2}{r}{\footnotesize Continua na próxima página} \\
\endfoot
\bottomrule
\endlastfoot
GET /health & Verificação de saúde da API. \\
GET /api/v1/events & Listagem paginada com filtros. \\
GET /api/v1/events/\{slug\} & Detalhamento de um evento. \\
GET /api/v1/admin/... & Indicadores administrativos e revisão. \\
\end{longtable}
\end{center}

\section{Filtros de consulta}

\begin{lstlisting}
q=<texto>
event_status=upcoming|ongoing|finished|active|all
submission_status=open|closed|unknown|all
registration_status=open|closed|unknown|all
area=<categoria>
region=latin-america|other|unknown
include_finished=true|false
limit=<número>
offset=<número>
\end{lstlisting}

\section{Exemplo de consulta regional}

\begin{lstlisting}
curl -s \
  'http://localhost:8000/api/v1/events?region=latin-america&include_finished=true&limit=30' | \
  docker compose exec -T api python -m json.tool
\end{lstlisting}

\section{Estrutura resumida da resposta}

\begin{lstlisting}
{
  "slug": "ichep-2026",
  "acronym": "ICHEP 2026",
  "event_status": "upcoming",
  "region": "latin-america",
  "calls": [...],
  "registrations": [...],
  "resources": [...],
  "confidence_score": 0.95,
  "confidence_level": "high",
  "verification_freshness": "current",
  "last_verified_at": "2026-07-26T..."
}
\end{lstlisting}
